\vspace{1.1em}

核心代码见 \texttt{cumcm/t4/}（全部实现），此处列出最关键的布局构造与"半圆盘命中"判据：

\begin{lstlisting}[language=Python]
# cumcm/t4/sweep.py —— 20 点定案布局
def measure_layout():
    """20 个测量位置：原点 + 7 覆盖基点 + 12 均匀方位外圈点（MID_RING_N=0 无内部补点）。"""
    base = np.asarray(SURVEY_CENTERS, dtype=float)      # 7 覆盖基点（复用问题三）
    ang = np.arange(OUTER_RING_N, dtype=float) * (2.0 * math.pi / OUTER_RING_N)
    ring = np.stack((OUTER_RING_RAD * np.cos(ang), OUTER_RING_RAD * np.sin(ang)), axis=1)
    mid_pts = np.empty((0, 2))
    if MID_RING_N > 0:                                  # 23 点版（3 内部补点）仅复现用
        mid = np.arange(MID_RING_N, dtype=float) * (2.0 * math.pi / MID_RING_N) \
            + math.radians(MID_RING_ROT_DEG)
        mid_pts = np.stack((MID_RING_RAD * np.cos(mid), MID_RING_RAD * np.sin(mid)), axis=1)
    return [(0.0, 0.0)] + [(float(x), float(y)) for x, y in base] + \
        [(float(x), float(y)) for x, y in mid_pts] + \
        [(float(x), float(y)) for x, y in ring]
\end{lstlisting}

\begin{lstlisting}[language=Python]
# cumcm/t4/sweep.py —— 定向源"半圆盘命中"判据（与官方引擎同构）
def _directional_hit(g, theta_deg, R, p):
    """测量点 p 能否听到位于 g、波束方向 theta_deg、半径 R 的定向源。"""
    if np.linalg.norm(p - g) > R:            # 距离超接收半径
        return False
    u = np.array([math.cos(math.radians(theta_deg)), math.sin(math.radians(theta_deg))])
    return float((p - g) @ u) >= 0.0         # 在波束方向前向半平面内（含边界）
\end{lstlisting}

完整可复现命令（见第 7 节复现说明）：

\begin{verbatim}
python T4.py --plan-only            # 扫描方案 + 听到率统计（不连模拟器）
python T4.py --practice 20 --seed 0 # 20 局演练（同 seed 逐字节可复现）
python T4.py                        # 官方模式：连 http://127.0.0.1:2026
python T4_figures.py                # 从 results/t4/ 生成论文图表
\end{verbatim}